\documentstyle[%


righttag%


,11pt%


,amssymb%


,russian%               remove in Eng.


,izvru%                 Set izven% in Eng.


]{amsart}





%





\newcommand{\la}{\langle}


\newcommand{\ra}{\rangle}


\newcommand{\up}[2]{\overset{#1}{#2}{}}


\newcommand{\diag}{\operatorname{diag}}


\newcommand{\tts}[1]{}


\numberwithin{equation}{section}





\theoremstyle{definition}


\theorembodyfont{\rm}


\newtheorem{defns}{\hskip\parindent Определение}[section]





%\hoffset=-15mm%                  For Eng.


\setcounter{page}{68}


\god{2001}


\nomer{12~(475)} %                        Remove in Eng.


%\nomer{Vol.\,45, No.\,12}%               Set in Eng.


%\str{\pageref{begin}--\pageref{end}}%  Set in Eng.


\udk{517.977.56}





\author{\it В.А.\,Терлецкий}


% NB:   ^^ remove in Eng.


\title{Обобщенное решение многомерных полулинейных гиперболических систем}





\date{}


%\pagestyle{myheadings}%         Set in Eng.


\pagestyle{plain} %              Remove in Eng.


\begin{document}


%\thispagestyle{plain}%          Set in Eng.


\maketitle


\label{begin}








При исследовании задач оптимального управления


разрывность управляющих воздействий создает необходимость рассматривать


решения дифференциальных уравнений, определяющих допустимый процесс, в


обобщенном смысле. Особенно остро эта проблема стоит для систем уравнений 


с частными производными, где зачастую невозможно построение не только 


гладкого, но и просто непрерывного решения, соответствующего допустимому 


управлению.





В основу понятия обобщенного решения могут быть положены самые различные 


подходы: интегральные законы сохранения, метод искусственной вязкости, способ


предельного перехода в разностных аппроксимациях, аппарат теории обобщенных


функций, понятие потенциала решения, а также другие схемы [1], [2].


К сожалению, далеко не все из них приемлемы для использования 


в теории оптимального управления. Например, в гиперболических системах 


многомерных уравнений одно из наиболее сильных понятий обобщенного решения


было введено Фридрихсом [3]. Им и его последователями доказано существование 


и единственность такого обобщенного решения, которое допускает сколь угодно 


точную аппроксимацию по норме пространства $L_2$ последовательностью гладких


приближенных решений [3]--[9]. Но так как доказательство проводится на основе


интегральных операторов осреднения, оценка решения имеет вид энергетического 


неравенства, а значит, не может гарантировать малость возмущения решения


в точках области независимых переменных при малых возмущениях неоднородной 


части системы и начально-краевых условий. Таким образом, способ 


конструирования обобщенного решения в рамках данного подхода с 


принципиальной точки зрения не позволяет не только обосновать широко 


известный в теории оптимального управления метод приращений, но и уловить 


качественное различие в поведении приращения траектории на различных типах 


игольчатых вариаций [10].





В работах [11]--[16] использовалось понятие обобщенного решения, лишенное 


отмеченных недостатков. Оно основано на следующей достаточно общей идее из


[1]: для системы дифференциальных уравнений нужно построить эквивалентную


ей интегральную систему, решение которой и объявить обобщенным. Естественное 


стремление распространить этот подход на гиперболические системы многомерных 


уравнений долгое время упиралось в построение соответствующего ей 


интегрального эквивалента. В статье предлагается решение этой проблемы на 


базе многомерного аналога инвариантов Римана.





\section{Постановка задачи}





В $m+1$-мерном пространстве независимых переменных $(s,t)$, 


$s = (s_1,s_2,\dotsc,s_m)$ рассмотрим параллелепипед


${\mathrm P} = S \times T$, 


$S = S^{(1)} \times S^{(2)} \times \dotsb \times S^{(m)}$, 


$S^{(j)} = [s_{0_j}, s_{1_j}]$, $j=1,2,\dotsc,m$, $T = [t_0,t_1]$.


Пусть $\Omega$, $G$, ${\cal D}_0$, ${\cal D}_1$ --- соответственно


его полная, боковая, нижняя и верхняя поверхности,


$\Omega = G \cup {\cal D}_0 \cup {\cal D}_1$.


В параллелепипеде $\mathrm P$ определим систему дифференциальных уравнений


с частными производными


\begin{equation}


x_t+\sum_{j=1}^{m} A_{j}(s,t)x_{s_j} = f(x,s,t).


\end{equation}


Здесь $x=x(s,t)$, $x(s,t) \in {\mathrm E}^n$ --- решение, 


$A_j=A_j(s,t)$, $A_j(s,t)\in{\mathrm E}^{n\times n}$, $j=1,2,\dotsc,m$, ---


заданные матричные функции, $f=f(x,s,t)$, $f(x,s,t) \in {\mathrm E}^n$ ---


заданная вектор-функция.


Предположим, что элементы матриц $A_j$, $j=1,2,\dotsc,m$, 


непрерывны и имеют непрерывную по Липшицу частную производную по $s$ в области 


$\mathrm P$. Кроме того, будет сформулировано еще одно предположение 


относительно матриц $A_j$, $j=1,2,\dotsc,m$, которое позволяло бы считать 


дифференциальный оператор


\begin{equation}


Dx=x_t+\sum_{j=1}^{m} A_{j}(s,t)x_{s_j}


\end{equation}


гиперболическим в области $\mathrm P$.


Как известно [1], в одномерном (т.\,е. при


$m=1$) случае для этого необходимо и достаточно, чтобы матрица при частной 


производной $x_s$ была простой. Подчеркнем, что указанное условие является


минимальным в том смысле, что другие определения гиперболических систем 


(напр., [2]--[9]) используют более жесткие ограничения: симметричность 


матрицы или наличие у нее $n$ различных всюду в $\mathrm P$ вещественных 


собственных значений. Известные автору теоремы существования и единственности


решения многомерных $(m \geq 2)$ систем вида (1.1) получены в предположении 


симметричности матриц $A_j$, $j=1,2,\dotsc,m$. Здесь будем использовать 





\begin{defns}


Систему (1.1) назовем гиперболической, если для


почти всех $\nu \in {\mathrm E}^m$ и всех $(s,t) \in {\mathrm P}$ матрица


$$


{\cal A}(\nu,s,t)=\sum_{j=1}^{m} \nu_j A_{j}(s,t)


$$


имеет вещественные (необязательно различные) собственные значения 


$\lambda_i=\lambda_i (\nu,s,t)$, $i=1,\dotsc,n$, и не является дефектной,


т.\,е.


алгебраическая и геометрическая кратность ее собственных значений одинакова.


\end{defns}





Для недефектной матрицы ${\cal A}={\cal A}(\nu,s,t)$ всегда существуют 


базисы пространства ${\mathrm E}^n$, составленные из ее левых 


$\ell^{(i)}=\ell^{(i)}(\nu,s,t)$ и правых $p^{(i)}=p^{(i)}(\nu,s,t)$,


$i=1,2,\dotsc,n$, собственных векторов. Указанные базисы всегда 


можно построить так, чтобы они были биортогональными [17], т.\,е.


существуют такие


матрицы ${\cal L}=(\ell^{(1)},\ell^{(2)},\dotsc,\ell^{(n)})$,


${\cal P}=(p^{(1)},p^{(2)},\dotsc,p^{(n)})$, что будут верны соотношения


\begin{equation}


\begin{array}{c}


{\cal L}^{\prime}(\nu,s,t){\cal A}(\nu,s,t)=


\Lambda(\nu,s,t) {\cal L}^{\prime}(\nu,s,t),\\[2mm]


{\cal A}(\nu,s,t){\cal P}(\nu,s,t)=


{\cal P}(\nu,s,t) \Lambda (\nu,s,t),\ \;


{\cal L}^{\prime}(\nu,s,t){\cal P} (\nu,s,t)=\dot{\mathrm I}.


\end{array}


\end{equation}


Здесь $\Lambda=\diag \{ \lambda_1, \lambda_2,\dotsc, \lambda_n\}$, 


$^\prime$ -- знак транспонирования, $\Dot{\mathrm I}$ --- единичная матрица.





Легко заметить, что матрица ${\cal A}$ может быть недефектной только в том 


случае, если все матрицы $A_{j}$, $j=1,2,\dotsc,m$, недефектные. Обратное 


утверждать, вообще говоря, нельзя.


Отметим, что предположение о симметричности матриц, формирующих 


дифференциальный оператор (1.2), равносильно [18] равенству 


${\cal L}={\cal P}$, т.\,е. приводит к частному случаю соотношений (1.3) и, 


следовательно, принятое здесь в определении 1.1


понятие гиперболичности системы заведомо не уже, чем те, 


которые используются в [4]--[9]. 





Будем считать вектор-функцию $f(x,s,t)$ непрерывной по Липшицу по переменной 


$x$ при фиксированных


$(s,t) \in {\mathrm P}$ и интегрируемой по Лебегу в $\mathrm P$ 


при фиксированных $x \in {\mathrm E}^n$.





При сделанных предположениях относительно дифференциального оператора (1.2) и


вектор-функции $f(x,s,t)$ требуется построить обобщенное решение системы (1.1),


обосновать его существование и единственность в случае корректного способа 


постановки смешанных условий для решения $x$ на границе $\Omega$ области


$\mathrm P$ и установить оценки роста решения почти всюду в $\mathrm P$.





\section{Характеристики и бихарактеристики гиперболического оператора}





Введем характеристики гиперболического оператора (1.2). Пусть


в $\mathrm P$ уравнением $\chi (s,t)=с$, $\chi \in C_1({\mathrm P})$, задана 


поверхность ${\cal X}_с$. Говорят, что поверхность ${\cal X}_с$ является 


характеристикой гиперболического оператора (1.2), если


$$


\det \bigg| \chi_t (s,t) \Dot{\mathrm I}+


\sum_{j=1}^{m} {\chi_s}_j A_{j}(s,t) \bigg|=0


$$


или если при некотором $i=1,2,\dotsc,n$ выполняется равенство


\begin{equation}


\chi_t (s,t)+\lambda_i (\chi_s,s,t)=0.


\end{equation}





Отметим, что собственные значения $\lambda_i (\nu,s,t)$ и собственные векторы


$\ell^{(i)}(\nu,s,t)$, $p^{(i)}(\nu,s,t)$ матрицы ${\cal A}(\nu,s,t)$ 


наследуют гладкость матриц $A_{j}$,


$j=1,2,\dotsc,m$, по $(s,t) \in {\mathrm P}$.


По переменной $\nu$ они обладают свойством однородности, в чем нетрудно 


убедиться, использовав линейность по $\nu$ матрицы ${\cal A}(\nu,s,t)$ и 


соотношения (1.3). При этом функции $\lambda_i (\nu,s,t)$ однородны первой 


степени, а вектор-функции $\ell^{(i)}(\nu,s,t)$, $p^{(i)}(\nu,s,t)$ --- 


нулевой степени, чего всегда можно добиться их нормировкой, 


ибо $\la\ell^{(i)}(\nu,s,t), p^{(i)}(\nu,s,t)\ra \equiv 1$,


$i=1,2,\dotsc,n$.


Итак, имеем


\begin{equation}


\begin{array}{c}


\lambda_i (\alpha \nu,s,t)=\alpha \lambda_i (\nu,s,t), \ \;


\ell^{(i)}(\alpha \nu,s,t)=\ell^{(i)}(\nu,s,t), \\[2mm]


p^{(i)}(\alpha \nu,s,t)=p^{(i)}(\nu,s,t),\ \; \alpha \geq 0,\ \;


i=1,2,\dotsc,n.


\end{array} 


\end{equation}





Свойство (2.2) позволяет применить эффективный метод построения решения 


нелинейного дифференциального уравнения (2.1) по аналогии с тем, как это 


сделано в [2] для двумерного $(m=2)$ случая.





Действительно, по теореме Эйлера для однородных первой степени функций 


справедливы тождества


\begin{equation}


\lambda_i (\nu,s,t)=\la {\lambda_i}_{\nu} (\nu,s,t), \nu \ra, \ \;


\nu \in {\mathrm E}^m,\ \; (s,t) \in {\mathrm P},\ \; i=1,2,\dotsc,n.


\end{equation}


Положив здесь для краткости $\nu=\chi_s$, перепишем (2.1) в виде


\begin{equation}


\chi_t+\la {\lambda_i}_{\nu} (\chi_s,s,t), \chi_s\ra=0, \ \;


i=1,2,\dotsc,n.


\end{equation}


Каждое из уравнений (2.4) может быть проинтегрировано методом характеристик 


([19], c.\,247) с незначительным его видоизменением, обусловленным следующим 


обстоятельством. При частных производных решения в (2.4) стоят коэффициенты, 


зависящие не от самого решения $\chi$, что позволило бы использовать метод


характеристик стандартным образом, а функции от тех же частных производных


решения. Поэтому характеристики уравнений (2.4) зависят непосредственно не от


решения $\chi$, а от его градиента по пространственным переменным $\chi_s$.


Следовательно, для нахождения решения системы обыкновенных дифференциальных


уравнений, устанавливающего траекторию характеристики уравнения из (2.4),


необходимо одновременно с ним уметь строить вектор-функцию $\chi_s$ вдоль этой 


характеристики.





Получим систему обыкновенных дифференциальных уравнений, решением которой и 


является градиент $\chi_s$ функции $\chi$. Для этого продифференцируем 


уравнение (2.1) по $s$ и воспользуемся обозначением $\nu=\chi_s$.


Будем иметь для каждого собственного значения $\lambda_i$ систему $2m$


обыкновенных дифференциальных уравнений


\begin{equation}


\frac{ds}{dt}={\lambda_i}_\nu (\nu,s,t), \quad


\frac{d\nu}{dt}=- {\lambda_i}_s (\nu,s,t),\quad i=1,2,\dotsc,n.


\end{equation}





В силу однородности нулевой степени по $\nu$ вектор-функций ${\lambda_i}_\nu$ 


и непрерывности по Липшицу по переменной $s$ вектор-функций ${\lambda_i}_s$


можно гарантировать, что для любого фиксированного 


$(\mu,\xi,\tau) \in {\mathrm E}^m \times S \times T$ существует единственное 


абсолютно непрерывное решение системы (2.5), которое будем называть, 


заимствуя термин из [2], бихарактеристикой и обозначать вектор-функциями 


$s^{(i)}(\mu,\xi,\tau;t)$ и ${\nu}^{(i)}(\mu,\xi,\tau;t)$ соответственно.





Понятно, что вдоль интегральной кривой $s=s^{(i)}(\mu,\xi,\tau;t)$ функция


$\chi_i$ сохраняет свое значение, т.\,к. с учетом (2.4)


$$


\frac{d}{dt} \chi_i (s^{(i)}(\mu,\xi,\tau;t),t)={\chi_i}_t+


\sum_{j=1}^{m} {\lambda_i}_{\nu_j} (\nu,s,t) {\chi_i}_{s_j}=0,\quad


i=1,2,\dotsc,n. 


$$


На этом свойстве основан следующий алгоритм построения конкретной 


характеристики ${\cal X}_с^{(i)}$. Во-первых, для уравнения (2.1) или 


эквивалентного ему уравнения (2.4) следует задать начальное условие. 


Пусть оно имеет вид 


\begin{equation}


\chi _i (s,\tau)=\zeta(s),


\end{equation}


где $\zeta=\zeta(s)$ --- некоторая кусочно-гладкая функция с невырожденным 


всюду в $S$ градиентом


$\| \nabla \zeta(s) \| > 0$. Во-вторых, необходимо найти 


точку $\xi \in S$, удовлетворяющую уравнению $\zeta (\xi)=c$, и вычислить


в ней градиент $\mu=\nabla\zeta(\xi)$. В третьих, нужно проинтегрировать 


систему (2.5) с условиями Коши $t=\tau$, $\nu=\mu$, $s=\xi$. Получим


бихарактеристику $s=s^{(i)}(\mu,\xi,\tau;t)$ и вектор-функцию 


$\nu={\nu}^{(i)}(\mu,\xi,\tau;t)$, которая, как уже отмечалось, совпадает с 


градиентом ${\chi_i}_s$ в точках


$(s^{(i)}(\mu,\xi,\tau;t), t) \in {\mathrm P}$.


Другими словами, вдоль бихарактеристики функция $\chi_i$ и ее градиент


${\chi_i}_s$ вычисляются по формулам $\chi_i (s^{(i)}(\mu,\xi,\tau;t), t)=


\zeta (s^{(i)}(\mu,\xi,\tau;t))$, ${\chi_i}_s (s^{(i)}(\mu,\xi,\tau;t), t)=


\nu^{(i)}(\mu,\xi,\tau;t)$.





Описанная схема дает эффективный инструмент для непосредственного построения 


$m$-мерной характеристики ${\cal X}_с^{(i)}$ по заданной в момент $t=\tau$ 


уравнением $\zeta(s)=c$ ``начальной'' $m{-}1$-мерной поверхности. Кроме того, 


она подсказывает структуру связи между функциями $\chi_i$ и $\zeta$. Покажем,


что с формальной точки зрения решение задачи (2.1), (2.6) определяет 


соотношение


\begin{equation}


{\chi_i} (\zeta(\cdot),\tau;s,t)=


\zeta(s^{(i)} ( \nu^{(i)}(\mu,\xi,\tau;t),s,t; \tau)).


\end{equation}





Справедливость начальных условий (2.6) для функции (2.7) очевидна из 


принятого при определении бихарактеристики свойства 


$s^{(i)}(\nu,s,\tau;\tau) \equiv s$. При проверке (2.1) для краткости 


обозначим $\nu=\nu^{(i)}(\mu,\xi,\tau;t)$, $\xi=s^{(i)}(\nu,s,t; \tau)$.


Непосредственно из (2.7) и второго равенства в (2.5) находим частные 


производные


\begin{gather*}


{\chi_i}_t=\la \nabla\zeta(\xi), s^{(i)}_t (\nu,s,t;\tau) -


s^{(i)}_\nu (\nu,s,t;\tau) {\lambda_i}_s (\nu,s,t) \ra, \\


{\chi_i}_s={s^{(i)}_s (\nu,s,t;\tau)}^{\prime} 


\nabla \zeta(\xi).


\end{gather*} 


После подстановки этих выражений в (2.4) будем иметь


$$


{\chi_i}_t+\la {\lambda_i}_\nu (\chi_s,s,t), \chi_s\ra=


\la \nabla\zeta(\xi), s^{(i)}_t (\nu,s,t;\tau)+


s^{(i)}_s (\nu,s,t;\tau) {\lambda_i}_\nu (\nu,s,t) -


s^{(i)}_\nu (\nu,s,t;\tau) {\lambda_i}_s (\nu,s,t) \ra.


$$


Второй сомножитель в этом скалярном произведении равен нулю в силу тождества


\begin{equation}


s^{(i)}_t (\nu,s,t;\tau)+s^{(i)}_s (\nu,s,t;\tau) {\lambda_i}_\nu (\nu,s,t) 


- s^{(i)}_\nu (\nu,s,t;\tau) {\lambda_i}_s (\nu,s,t) \equiv 0,


\end{equation}


вытекающего из тавтологии 


$$


\xi \equiv s^{(i)}(\nu^{(i)}(\mu,\xi,\tau;t), 


s^{(i)} (\mu,\xi,\tau;t),t;\tau),\ \; \mu \in {\mathrm E}^m, 


\ \; (\xi,\tau) \in {\mathrm P},\ \; t \in T.


$$


Таким образом, функция (2.7) есть решение задачи Коши (2.1), (2.6). 


Следовательно, проблемы интегрирования нелинейного дифференциального 


уравнения (2.1) и построения характеристик дифференциального оператора (1.2)


могут решаться путем нахождения соответствующих решений бихарактеристической 


системы (2.5). Этим, как будет видно из дальнейшего, роль бихарактеристик


в построении обобщенного решения системы (1.1) далеко не ограничивается.





\section{Инварианты Римана}





По аналогии с тем, как это принято в системах одномерных гиперболических 


уравнений [1], [11], [16],


введем новые переменные $r=r(\nu,s,t)$,


$r(\nu,s,t) \in {\mathrm E}^n$,


$(\nu,s,t) \in {\mathrm E}^m \times S \times T$,


связав их с решением $x$ системы (1.1) линейным взаимно однозначным 


преобразованием 


\begin{equation}


r(\nu,s,t)={\cal L}^{\prime}(\nu,s,t) x(s,t),\quad


x(s,t)={\cal P}(\nu,s,t) r(\nu,s,t).


\end{equation}


Следуя традиции, функции $r_i=r_i(\nu,s,t)$, $i=1,2,\dotsc,n$, будем


по-прежнему называть инвариантами Римана, несмотря на их существенное 


отличие от одномерного прототипа --- наличие наряду с пространственной $s$ 


и временной $t$ еще одной независимой переменной $\nu$, которая ассоциируется


с некоторым направлением в точке $s$.





Нашей ближайшей целью будет вывод системы дифференциальных уравнений, которой


должны подчиняться инварианты $r$ в случае, если решение $x$ системы (1.1)


является классическим, т.е. непрерывным и непрерывно дифференцируемым в 


области $\mathrm P$.





Вначале воспользуемся условиями (1.3) и проанализируем структуру матриц 


$$


M_j (\nu,s,t)={\cal L}^{\prime}(\nu,s,t) A_j(s,t){\cal P}(\nu,s,t),


\quad j=1,2,\dotsc,m.


$$


Опуская для краткости аргументы $\nu$, $s$, $t$,


заметим, что справедливы равенства


$$


\la \ell^{(i)}, {\cal A} p^{(i)} \ra=\lambda_i,\ \; i=1,2,\dotsc,n.


$$


Продифференцируем каждое из них по переменной $\nu_j$, $j=1,2,\dotsc,m$.


Будем иметь


$$


\lambda_{i_{\nu_j}}=\la {\ell^{(i)}}_{\nu_j}, {\cal A} p^{(i)} \ra+


\la\ell^{(i)}, {\cal A} {p^{(i)}}_{\nu_j}\ra+\la\ell^{(i)}, A_j p^{(i)}\ra=


{\lambda_i} (\la {\ell^{(i)}}_{\nu_j}, p^{(i)} \ra+


\la\ell^{(i)}, {p^{(i)}}_{\nu_j}\ra)+\la\ell^{(i)}, A_j p^{(i)} \ra.


$$


Так как $\la {\ell^{(i)}},p^{(i)} \ra \equiv 1$, то $\lambda_{i_{\nu_j}}=


\la\ell^{(i)}, A_j p^{(i)}\ra$ или в матричной форме


$\Lambda_{\nu_j}=\diag \{ M^{(1,1)}_j,$ $M^{(2,2)}_j,\dotsc, M^{(n,n)}_j\}$.


%%%%%%%%% join in eng.               ^^^^^^


Пусть $\up{\vee}{M_j}=M_j - \Lambda_{\nu_j}$, $j=1,2,\dotsc,m$.


Вообще говоря, эти матричные функции нетривиальны, т.\,к. одновременная 


диагонализация всех матриц $A_1,A_2,\dotsc,A_m$ одним линейным преобразованием,


как известно [18], возможна лишь для коммутативного относительно операции 


умножения семейства матриц, т.\,е. для случая $A_j A_k=A_k A_j$, 


$j,k=1,2,\dotsc,m$. Такое предположение является очень жестким, оно резко 


сужает возможные приложения системы (1.1) и, естественно, упрощает проблему 


построения обобщенного решения фактически до ее уровня в одномерном $(m=1)$ 


варианте. Не вдаваясь в подробности, отметим лишь, что у 


дифференциального оператора (1.2) с коммутативными матрицами 


$A_1,A_2,\dotsc,A_m$ характеристики совпадают с бихарактеристиками


и не зависят от $\nu$.


Тем не менее для совокупности матриц $\up{\vee}{M_j}$, $j=1,2,\dotsc,m$,


как нетрудно видеть, выполняется тождество


\begin{equation}


\sum_{j=1}^m \nu_j \up{\vee}{M_j}(\nu,s,t) \equiv 0, 


\ \; \nu \in {\mathrm E}^m,\ \; (s,t) \in {\mathrm P}.


\end{equation}





Подставим теперь в систему (1.1) вместо решения $x$ соответствующее ему


представление из (3.1). Получим


\begin{multline*}


Dx=x_t+\sum_{j=1}^m A_j x_{s_j}=({\cal P} r)_t+


\sum_{j=1}^m A_j ({\cal P} r)_{s_j}=


D{\cal P} \cdot r+{\cal P} (r_t+{\cal L}^{\prime}


\sum_{j=1}^m A_j {{\cal P} r}_{s_j})= \\


%


=D{\cal P} \cdot r+{\cal P} (r_t+\sum_{j=1}^m \Lambda_{\nu_j} r_{s_j}+


\sum_{j=1}^m \up{\vee}{M_j} r_{s_j})=f({\cal P}r,s,t).


\end{multline*}


Отсюда следует


\begin{equation}


r_t+\sum_{j=1}^m \Lambda_{\nu_j} r_{s_j}={\cal L}^{\prime} 


(f - D {\cal P} \cdot r) - \sum_{j=1}^m \up{\vee}{M_j} r_{s_j}.


\end{equation}


Подсчитаем вдоль произвольной бихарактеристики 


$\xi=s^{(i)}(\nu,s,t;\tau)$, $\mu=\nu^{(i)}(\nu,s,t;\tau)$ полную 


производную $i$-го инварианта Римана 


$r_i=r_i ( \nu^{(i)}(\nu,s,t;\tau), s^{(i)}(\nu,s,t;\tau),\tau)$ по


переменной $\tau$


\begin{equation}


\frac{d}{d\tau}r_i=r_{i_\tau}+\sum_{j=1}^m {\lambda_i}_{\mu_j} 


{r_i}_{\xi_j} - \sum_{j=1}^m {\lambda_i}_{\xi_j} {r_i}_{\mu_j},


\ \; i=1,2,\dotsc,n.


\end{equation}


Сравнив это выражение с левой частью $i$-го уравнения системы (3.3), заметим,


что она отличается от полной производной инварианта $r_i$ слагаемым 


$-\la\lambda_{i_s}, r_{i_\nu}\ra$. Поэтому после суммирования соответствующих 


уравнений системы (3.3) с равенствами


$$


- \la\lambda_{i_s}, r_{i_\nu}\ra=


- \la\lambda_{i_s},\ell^{(i)\prime}_\nu {\cal P} r\ra,\ \; i=1,2,\dotsc,n,


$$


которые справедливы в силу соотношений (3.1), получим нужный вид 


инвариантной системы


\begin{multline}


{r_i}_t+\la\lambda_{i_\nu}, r_{i_s}\ra - \la\lambda_{i_s}, r_{i_\nu}\ra=


\la\ell^{(i)}, f - D{\cal P} \cdot r\ra - 


\la\lambda_{i_s},\ell^{(i)\prime}_\nu {\cal P} r\ra -\\


%


- \sum_{j=1}^m \la \up{\vee}{M}_{j}^{(i)}, r_{s_j}\ra, \quad


i=1,2,\dotsc,n,\ \; \nu \in {\mathrm E}^m,\ \; (s,t) \in {\mathrm P}.


\end{multline}


Здесь $\up{\vee}{M}_j^{(i)}$ --- $i$-я строка матрицы 


$\up{\vee}{M_j}$.


Пусть бихарактеристика $\xi=s^{(i)}(\nu,s,t;\tau)$, выпущенная из точки


$s$ в момент времени $t$ по направлению вектора $\nu$ в ``обратном'' времени


$\tau<t$, достигает границы $\Omega$ области $\mathrm P$ в точке


$\xi^{(i)}$ в момент времени $\tau^{(i)}$ с направлением $\mu^{(i)}$.


Понятно, что для построения такой тройки


$\mu^{(i)}$, $\xi^{(i)}$, $\tau^{(i)}$


достаточно проинтегрировать $i$-ю систему обыкновенных дифференциальных


уравнений (2.5) из точки $(\nu,s,t)$ по $\tau<t$. Как уже 


отмечалось, левая часть каждого из уравнений системы (3.5) представляет собой 


полную производную (3.4) соответствующего инварианта Римана. Следовательно,


решение дифференциальной системы (3.5) должно удовлетворять интегральной 


системе


\begin{multline}


r_i(\nu,s,t)=r_i(\mu^{(i)}, \xi^{(i)}, \tau^{(i)})+


\int_{\tau^{(i)}}^{t} \bigg[ \la\ell^{(i)}, f - D{\cal P}\cdot r\ra -


\la\lambda_{i_s},\ell^{(i)\prime}_\nu {\cal P} r\ra - \\


%


- \sum_{j=1}^m \la \up{\vee}{M}_j^{(i)}, r_{s_j}\ra \bigg]


\bigg|\begin{Sb}


\xi=s^{(i)}(\nu,s,t;\tau) \\


\mu=\nu^{(i)}(\nu,s,t;\tau)


\end{Sb} 


d\tau,\ \; i=1,2,\dotsc,n,\ \; \nu \in {\mathrm E}^m, 


\ \; (s,t) \in {\mathrm P}.


\end{multline}


С помощью непосредственной громоздкой проверки, которую


приводить не будем, убедимся в справедливости обратного 


утверждения. Отметим только, что при этом существенным образом наряду с (2.8)


используется симметричное ему тождество


$$


\nu^{(i)}_t (\nu,s,t;\tau)+\nu^{(i)}_s (\nu,s,t;\tau) 


{\lambda_i}_\nu (\nu,s,t) - \nu^{(i)}_\nu (\nu,s,t;\tau) 


{\lambda_i}_s (\nu,s,t) \equiv 0,


$$


которое тоже вытекает из тавтологии 


$$


\mu \equiv \nu^{(i)}(\nu^{(i)}(\mu,\xi,\tau;t), 


s^{(i)} (\mu,\xi,\tau;t),t;\tau),\ \; \mu \in {\mathrm E}^m, 


\ \; (\xi,\tau) \in {\mathrm P},\ \; t \in T.


$$


Таким образом, при условии существования классического решения $x$ системы


(1.1) эквивалентность дифференциальной (3.5) и интегральной (3.6) систем 


доказана.





\section{Построение корректных смешанных условий}





Смешанные условия для системы (1.1) сконструируем по следующей схеме.


Вначале выясним, каким образом можно определить такие условия для инвариантной


системы (3.5), а затем трансформируем их на систему (1.1), пользуясь 


взаимно однозначным отображением (3.1).





Из (3.6) видно, что для построения единственного решения инвариантной 


системы (3.5) нужно его ``закрепить'' в начальных точках 


$(\mu^{(i)},\xi^{(i)},\tau^{(i)})$ каждой бихарактеристики 


$\xi=s^{(i)}(\nu,s,t;\tau)$, $\nu \in {\mathrm E}^m$,


$(s,t) \in {\mathrm P}$,


как если бы речь шла о постановке условий Коши для бесконечномерной системы 


обыкновенных дифференциальных уравнений (3.5).





Введем в рассмотрение функции 


$\sigma_i{=}\sigma_i(\mu,\xi,\tau)$, $i{=}1,2,\dotsc,n$, положив 


$\sigma_i{=}\nu_0+\la{\lambda_i}_\nu (\mu,\xi,\tau), \nu\ra$, где 


$\wh\nu=(\nu_0,\nu)$ --- вектор внешней нормали к поверхности $\Omega$ 


области $\mathrm P$ в точке $(\xi, \tau)$, а $\mu$ --- произвольный вектор из


${\mathrm E}^m$. С помощью этих функций удобно классифицировать произвольную 


тройку $(\mu,\xi,\tau)$, $\mu \in {\mathrm E}^m$, $(\xi,\tau) \in \Omega$, 


как начальную или конечную точку конкретной бихарактеристики. Будем говорить,


что такая тройка является начальной (конечной) точкой бихарактеристики 


$s=s^{(i)}(\mu,\xi,\tau;t)$, если $\sigma_i < 0$ ($\sigma_i > 0$).


Введенное определение представляется вполне естественным аналитическим 


критерием проверки геометрического факта входа в область ${\mathrm P}$ или


выхода из нее в точке $(\xi,\tau)$ гладкой кривой с касательным вектором


$(ds,dt)=({\lambda_i}_\nu (\mu,\xi,\tau), 1)dt$. Случай $\sigma_i=0$


соответствует ``гладкому прикосновению''


бихарактеристики к границе области и не позволяет 


на основе используемой выше локальной информации идентифицировать


точку $(\mu,\xi,\tau)$ более детально. Вообще говоря, она может быть и 


начальной, и конечной, и промежуточной.


Вместе с тем, очевидная ``экзотичность''


равенства $\sigma_i=0$ предопределяет нулевую меру множества тех точек из


$\mathrm P$, для которых оно влияет на решение гиперболической 


системы.





Обратимся теперь к рассмотрению механизма смешанных условий, зафиксировав


для простоты некоторую точку $(\xi, \tau) \in \Omega$. Знак функции 


$\sigma_i=\sigma_i (\mu, \xi,\tau)$ служит индикатором всей тройки 


$(\mu,\xi,\tau)$ в том смысле, что либо требует ``закрепить'' значение 


инварианта Римана $r_i (\mu,\xi,\tau)$ ``извне'' при $\sigma_i < 0$, либо,


если $\sigma_i \geq 0$, запрещает такое закрепление, т.\,к. в этом 


случае значение $r_i (\mu,\xi,\tau)$ определяется при решении дифференциальной


системы в области $\mathrm P$ с использованием условия Коши


в точке, соответствующей 


данной бихарактеристике. Предположим, что $\sigma_i < 0$ при выбранном 


$\mu \in {\mathrm E}^m$. Значит, для $r_i (\mu,\xi,\tau)$ 


начальное условие ставить нужно, хотя оно не может быть совсем


произвольным, ибо в силу (3.1) тождество 


$r_i (\mu,\xi,\tau)=\la\ell^{(i)} (\mu,\xi,\tau), x(\xi,\tau)\ra$ должно 


сохраняться для всех $\mu \in {\mathrm E}^m$, на которых функция 


$\sigma_i (\mu,\xi,\tau) < 0$. Другими словами, следует обеспечить корректное


взаимодействие между различными проекциями $r_i (\mu,\xi,\tau)$ одного и того


же вектора $x(\xi,\tau)$ на направления $\ell^{(i)} (\mu,\xi,\tau)$.





Заметим, что среди всевозможных векторов $\mu \in {\mathrm E}^m$ особую роль 


играет подвектор $\nu$ вектора внешней нормали $\wh \nu$ к поверхности


$\Omega$ в точке $(\xi,\tau)$, т.\,к. лишь с его помощью можно выяснить,


приходит в точку $(\xi,\tau)$ (т.\,е. находится внутри $\mathrm P$ при 


$t < \tau$) или уходит из нее (т.\,е. находится внутри $\mathrm P$ при 


$t > \tau$) произвольная бихарактеристика 


$s=s^{(i)}(\mu,\xi,\tau;t)$, $\mu \in {\mathrm E}^m$. Кроме того, функции 


$\sigma_i$ вычисляются гораздо проще при условии $ \mu=\nu(\xi, \tau)$.


Действительно, в силу (2.3) имеем 


$\sigma_i (\nu,\xi,\tau)=\nu_0+{\lambda_i} (\nu,\xi,\tau)$. Приведенные


соображения говорят в пользу выбора в качестве ``базового'' направления в 


произвольной точке $(\xi,\tau) \in \Omega$ именно вектора $\nu$.





Определим вспомогательные конструкции, существенно упрощающие 


запись начально-кра\-е\-вых условий. Вначале построим матрицу


$$


B(\xi,\tau)=\nu_0(\xi,\tau){\mathrm I}+\sum_{j=1}^{m} \nu_j (\xi, \tau) 


A_{j}(\xi,\tau).


$$


С учетом соотношений (1.3) для нее справедливо также равенство 


$$


B(\xi,\tau)={\cal P} (\nu(\xi,\tau),\xi,\tau) \Sigma(\xi,\tau)


{\cal L}^{\prime} (\nu(\xi, \tau), \xi,\tau),


$$ 


в котором $\Sigma=\diag \{ \sigma_1,\sigma_2,\dotsc,\sigma_n \}$. Через


$\Sigma^{-}$ обозначим матрицу, получаемую из матрицы $\Sigma$


заменой положительных диагональных значений на нули. Тогда начально-граничные


условия для инвариантного вектора $r(\nu,\xi,\tau)$ можно записать в виде


\begin{equation}


\Sigma^{-} (\xi,\tau) [r(\nu,\xi,\tau) - r^0 (\xi,\tau)]=0.


\end{equation}


Здесь $r^0=r^0 (\xi,\tau)$ считаем заданной функцией. Ясно, что равенство


(4.1) эквивалентно соотношениям


\begin{equation}


r_i (\nu,\xi,\tau)=r^0_i (\xi,\tau), \sigma_i (\nu, \xi,\tau) < 0, 


\ \; i=1,2,\dotsc,n.


\end{equation}


Теперь с помощью линейного невырожденного преобразования (3.1) нетрудно 


переписать условия (4.1) в терминах переменных $x$. Положив 


$x^0 (\xi,\tau)={\cal P}(\nu(\xi,\tau),\xi,\tau) r^0 (\xi, \tau)$ и


$B^{-}(\xi,\tau)={\cal P}(\nu(\xi,\tau),\xi,\tau) 


\Sigma^{-}(\xi,\tau) {\cal L}^{\prime} (\nu(\xi, \tau), \xi,\tau)$,


будем иметь


\begin{equation}


B^{-}(\xi,\tau) [x(\xi, \tau) - x^{0} (\xi,\tau)]=0.


\end{equation}


Понятно, что условия (4.3) в отличие от условий (4.1) ввиду недиагональности


матрицы $B^{-}$ не допускают покомпонентную расшифровку типа (4.2).





Покажем теперь, как следует определять, не противоречa 


тождеству (3.1), начально-гра\-нич\-ные условия для инвариантов 


$r_i (\mu,\xi,\tau)$, если $\sigma_i (\mu, \xi,\tau) < 0$ и $\mu \neq \nu$.


Здесь, разумеется, уже нельзя использовать компоненты $r_i^0$ вектор-функции


$r^0$, соответствующие неравенствам $\sigma_i (\mu, \xi,\tau) \geq 0$, 


которые автоматически игнорируются условиями (4.1)--(4.3). Если 


$\sigma_i (\nu, \xi,\tau) \geq 0$, то значение $r_i (\nu,\xi,\tau)$ 


находится как результат интегрирования системы (3.5). Добавив к этим 


компонентам условия (4.2), будем иметь в распоряжении весь вектор 


$r(\nu,\xi,\tau)$, что позволяет подсчитать значение $x(\xi,\tau)$ по формуле


(3.1). Окончательно имеем: если $\sigma_i (\mu,\xi,\tau) < 0$, то


$r_i (\mu,\xi,\tau)=\la\ell^{(i)} (\mu,\xi,\tau), x(\xi,\tau)\ra$.


Таким образом,


полностью сформировать граничные условия для системы (3.5) в момент времени


$\tau$ можно только после того, как эта система проинтегрирована в полосе


$[t_0, \tau]$.





\section{Существование и единственность обобщенного решения} 





Как уже отмечалось, 


дифференциальная система (3.5) и интегральная система (3.6) на гладких 


решениях эквивалентны. В то же время, в отличие от дифференциальной, 


интегральная система не теряет смысл и в том случае, когда входные параметры


задачи --- вектор-функции $f$ и $r^0$ --- разрывны по независимым переменным.


Как известно [1], именно это обстоятельство позволяет определить обобщенное 


решение исходной дифференциальной системы как решение интегральной. 


Для одномерных инвариантных систем такое построение обобщенного решения 


использовалось в [11], [16]. Принципиальным отличием


рассматриваемой здесь интегральной системы для 


инвариантов Римана от аналогичной системы из [11], [16] является последнее 


слагаемоe в (3.6), содержащее частные производные по пространственным 


переменным от недиагональных компонент решения. Однако можно показать, 


что почти всюду в $\mathrm P$ величина евклидовых норм этих производных и 


решения $r$ имеет одинаковый порядок. После установления такой оценки 


доказательство существования и единственности решения интегральной 


системы (3.6) по существу проходит по схеме [11], использующей метод 


последовательных приближений и принцип интегрального сжатия.





Cуществование и единственность решения инвариантной 


системы наряду с наличием взаимно однозначного преобразования


(3.1) адекватно построению обобщенного решения исходной системы (1.1)


с начально-граничными условиями (4.3). Укажем теперь оценку решения $x$ 


в точках области $\mathrm P$ через входные параметры $f$ и $x^0$


начально-краевой задачи (1.1), (4.3).





Для произвольной фиксированной точки $(s,t) \in {\mathrm P}$ рассмотрим 


характеристики ${\cal X}^{(i)}={\cal X}^{(i)}(s,t)$, положив


$$


{\cal X}^{(i)}=\{ (\xi,\tau) \in {\mathrm P},\ \xi=s^{(i)}(\nu,s,t; \tau), 


\ \nu \in {\mathrm E}^m,\ \| \nu \|=1\},\ \; i=1,2,\dotsc,n.


$$


Заметим, что в силу равенств (3.2) вектор $\mu=\nu^{(i)}(\nu,s,t; \tau)$


всегда ортогонален векторам $(M^{(i,k)}_1, M^{(i,k)}_2,\dotsc, M^{(i,k)}_m)$,


$i,k=1,2,\dotsc,n$, $i \neq k$, а значит, в системах (3.5), (3.6) градиенты


по $s$ инвариантов Римана $r_k$ направлены по касательной к поверхностям 


${\cal X}^{(i)}$. Основываясь на этом факте и тождествах (3.1), можно 


получить оценку


$$


\|x(s,t)\| \leq K \bigg[ \sum_{i=1}^n \bigg(\;


\int\limits_{\Omega \cap {\cal X}^{(i)}} \| x^0 \| d\omega+


\int\limits_{{\cal X}^{(i)}} \| f \| d\omega \biggr)+


\int\limits_{\Omega \cap \partial \triangle (s,t)}


\| x^0 \| d\omega+


\iint\limits_{\triangle(s,t)} 


\| f \| d\xi d\tau \bigg],


$$


которая справедлива почти всюду в $\mathrm P$.


Здесь $K>0$ --- некоторая константа; 


$\triangle (s,t)$ --- область зависимости решения $x$ в точке 


$(s,t) \in {\mathrm P}$,


являющаяся объединением всех характеристических конусов, 


охватываемых поверхностями 


${\cal X}^{(i)}(s,t)$, $i=1,2,\dotsc,n$; $\partial \triangle(s,t)$ --- граница 


области $\triangle (s,t)$.





\begin{thebibliography}{99}





\bibitem{1}


Рождественский Б.Л., Яненко Н.А. {\it Системы квазилинейных уравнений


и их приложения к газовой динамике}. -- М.: Наука, 1978. -- 687 с.


\bibitem{2}


Годунов С.К. {\it Уравнения математической физики}. -- М.: Наука, 1979. -- 


392 с.


\bibitem{3}


Friedrichs K.O. {\it The identity of weak and strong extensions of


differential operator} // Trans. Amer. Math. Soc. -- 1944. -- V.\,55. --


P.\,132--151.


\bibitem{4}


Friedrichs K.O. {\it Symmetric hyperbolic linear differential equations}


// Comm. Pure Appl. Math. -- 1954. -- V.\,7. -- \No\,2. -- P.\,345--392.


\bibitem{5}


Friedrichs K.O. {\it Symmetric positive linear differential equations}


// Comm. Pure Appl. Math. -- 1958. -- V.\,11. -- P.\,333--418.


\bibitem{6}


Lax P.D., Phillips R.S. {\it Local boundary conditions for dissipative 


linear differential operators} // Comm. Pure Appl. Math. -- 1960. --


V.\,13. -- P.\,427--455.


\bibitem{7}


Sarason L. {\it On weak and strong solutions of boundary value problems}


// Comm. Pure Appl. Math. -- 1962. -- V.\,15. -- P.\,237--288.


\bibitem{8}


Sarason L. {\it On boundary value problems for hyperbolic equations}


// Comm. Pure Appl. Math. -- 1962. -- V.\,15. -- P.\,373--395.


\bibitem{9}


Phillips R.S., Sarason L. {\it Singular symmetric first order


differential operators} // J. Math. and Mech. -- 1966. -- V.\,15.


-- \No\,2. -- P.\,235--271.


\bibitem{10}


Терлецкий В.А. {\it Вариационный принцип максимума в управляемых


системах многомерных гиперболических уравнений} // Оптимизация, управление,


интеллект. -- Иркутск, 2000. -- \No\,4. -- C.\,93--110.


\bibitem{11} 


Терлецкий В.А. {\it Необходимые условия оптимальности и численные методы


оптимизации в системах полулинейных гиперболических уравнений}: Автореф.


дис. $\dotsc$ канд. физ.-матем. наук. -- Л., 1983. -- 17 с.


\bibitem{12}


Васильев О.В., Срочко В.А., Терлецкий В.А. {\it Методы оптимизации и


их приложения}. Ч.\,2. {\it Оптимальное управление}. -- Новосибирск: Наука,


1990. -- 151 с.


\bibitem{13}


Vasiliev O.V., Terletsky V.A., Arguchintsev A.V. {\it Iterative


processes in optimization of semilinear hyperbolic system} //


11-th IFAC World Congress. -- Tallinn, 1990. -- V.\,6. -- P.\,216--220.


\bibitem{14}


Васильев О.В., Терлецкий В.А. {\it Итерационные процессы решения задач


оптимального управления в системах с сосредоточенными и распределенными


параметрами, основанные на принципе максимума Л.С.\,Понтрягина} //


Оптимизация: Модели. Методы. Решения. -- Новосибирск: Наука, 1992. --


C.\,35--54.


\bibitem{15}


Васильев О.В., Терлецкий В.А., Болдонов А.В. {\it К исследованию


некоторых задач оптимального управления, возникающих в обратной проблеме


цунами} // Методы численного анализа и оптимизации. -- Новосибирск:


Наука, 1987. -- C.\,3--33.


\bibitem{16}


Терлецкий В.А. {\it Вариационный принцип максимума в управляемых системах


одномерных гиперболических уравнений} // Изв. вузов. Математика.


-- 1999. -- \No\,12. -- C.82--90.


\bibitem{17}


Ланкастер П. {\it Теория матриц}. -- М.: Наука, 1982. -- 269 с.


\bibitem{18}


Хорн Р., Джонсон Ч. {\it Матричный анализ}. -- М.: Мир, 1989. -- 625 с.


\bibitem{19}


Петровский И.Г. {\it Лекции по теории обыкновенных дифференциальных 


уравнений}. -- М.: Наука, 1964 г. -- 272 с.





\end{thebibliography}





\vskip 2cm





\post{\it Иркутский государственный университет}{\it Поступила}


\post{}{24.08.2001}





\label{end}


\end{document}


